\subsection{Case Study}
\label{sec:casestudy}

To examine how admission and pacing interact under queue pressure, we selected a run on S26 Ultra with lower optional-stage execution and more frequent pacing than the peer medians (Table~\ref{tab:casestudy_selection}).
The run used production-default 12MP under normal memory conditions and started at the guard's level-4 cutoff.
Figure~\ref{fig:casestudy_12mp} traces the run over 30 captures.

Our key observation is that pacing responded to rising queue pressure before admission began skipping optional stages.
Pacing began at capture 5, and backlog growth slowed while it was active.
Slack nevertheless narrowed to 2.1\% by capture 8; admission began skipping \(M\) at capture 9, after which backlog declined.
With optional work reduced, pacing was inactive over captures 13--20, and slack recovered to 27.3\% by capture 15.
When pressure rose again, pacing resumed at capture 21 before \(S\) skips began at capture 23; backlog then declined, and slack recovered to 45.6\% by capture 30.
This case study demonstrates that pacing and admission can jointly recover timeout slack by moderating queue growth early and reducing optional work as budgets tighten.

% Notes: docs/exhibits.md#tab_casestudy_selection
\begin{table}[t]
  \centering
  \setlength{\abovecaptionskip}{2pt}%
  \setlength{\belowcaptionskip}{3pt}%
  \caption{Case-study run compared with its peers.}
  \label{tab:casestudy_selection}

  {\scriptsize
  \renewcommand{\arraystretch}{1.15}%
  \setlength{\tabcolsep}{5pt}%

  \begin{tabular}{>{\raggedright\arraybackslash}p{0.376\columnwidth}
                  >{\raggedleft\arraybackslash}p{0.075\columnwidth}
                  >{\raggedleft\arraybackslash}p{0.075\columnwidth}
                  >{\centering\arraybackslash}p{0.150\columnwidth}
                  >{\centering\arraybackslash}p{0.075\columnwidth}}
    \toprule
    \makecell[bl]{\textbf{Metric}} &
    \makecell[br]{\textbf{This}\\\textbf{run}} &
    \makecell[br]{\textbf{Peer}\\\textbf{med.}} &
    \makecell[bc]{\textbf{Peer}\\\textbf{[min--max]}} &
    \makecell[bc]{\textbf{vs.}\\\textbf{med.}} \\
    \midrule
    \(M\)~executed \(\uparrow\)
      & 8 & 14 & [8--30] & worse \\
    \(S\)~executed \(\uparrow\)
      & 22 & 30 & [22--30] & worse \\
    Pacing activated (\%)\,\(\downarrow\)
      & 55.2 & 31.0 & [13.8--69.0] & worse \\
    Cumulative applied delay (s)\,\(\downarrow\)
      & 5.21 & 4.17 & [3.51--9.70] & worse \\
    \bottomrule
  \end{tabular}}
\end{table}

% Notes: docs/exhibits.md#fig_casestudy_12mp
\begin{figure}[t]
  \centering
  \setlength{\abovecaptionskip}{2pt}%
  \setlength{\belowcaptionskip}{2pt}%
  \providecommand{\caseStudyDataDir}{data/case_study}
  \providecommand{\csTimeoutRange}[2]{%
    \fill[cs tmo range] ({axis cs:#1,0}|-{rel axis cs:0,0})
      rectangle ({axis cs:#2,0}|-{rel axis cs:0,1});}
  \tikzset{
    cs tmo range/.style={red!70!black,opacity=0.10},
  }
  \pgfplotsset{
    cs executed/.style={
      only marks,mark=square*,mark size=1.35pt,
      draw=black,fill=black!72,line width=0.4pt,
    },
    cs skipped/.style={
      only marks,mark=square,mark size=1.35pt,
      draw=black!45,line width=0.4pt,
    },
    cs delay/.style={
      ybar,bar width=2.1pt,
      draw=blue!55!black,fill=blue!55!black,fill opacity=0.55,line width=0.35pt,
    },
    cs backlog/.style={
      line width=0.7pt,draw=blue!55!black,
      mark=*,mark size=0.9pt,
      mark options={solid,draw=none,fill=blue!55!black},
    },
    cs slack/.style={
      line width=0.7pt,draw=black!80,
      mark=*,mark size=0.9pt,
      mark options={solid,draw=none,fill=black!80},
    },
    cs floor/.style={red!70!black,line width=0.6pt,densely dashed},
    cs step/.style={
      const plot,no marks,line width=0.65pt,draw=black!80,
    },
  }

  \begin{tikzpicture}
    \begin{groupplot}[
      group style={
        group size=1 by 4,
        vertical sep=3.5pt,
        xticklabels at=edge bottom,
      },
      scale only axis=true,
      width=0.745\columnwidth,
      xmin=0.5,xmax=30.5,
      xtick={5,10,15,20,25,30},
      tick label style={font=\scriptsize},
      label style={font=\scriptsize},
      grid=major,
      major grid style={gray!16,line width=0.2pt},
      axis line style={line width=0.4pt},
      tick style={line width=0.4pt},
      tick align=outside,
      tick pos=left,
      ylabel style={
        at={(axis description cs:-0.15,0.5)},
        anchor=south,
        align=center,
      },
    ]

      \nextgroupplot[
        height=0.115\columnwidth,
        ymin=0.4,ymax=2.6,
        ytick={1,2},
        yticklabels={{\(S\)},{\(M\)}},
        yticklabel style={font=\tiny},
        grid=none,
        clip=false,
      ]
      \csTimeoutRange{8}{12}
      \addplot[cs executed] table[x=shot,y=lane,col sep=comma]
        {\caseStudyDataDir/12mp_normal_bokeh_exec.csv};
      \addplot[cs skipped] table[x=shot,y=lane,col sep=comma]
        {\caseStudyDataDir/12mp_normal_bokeh_skip.csv};
      \addplot[cs executed] table[x=shot,y=lane,col sep=comma]
        {\caseStudyDataDir/12mp_normal_filter_exec.csv};
      \addplot[cs skipped] table[x=shot,y=lane,col sep=comma]
        {\caseStudyDataDir/12mp_normal_filter_skip.csv};
      \node[anchor=south,font=\tiny,align=center,
            text=red!70!black,inner sep=1pt]
        at ([yshift=1pt]axis cs:10,2.6)
        {baseline\\first-timeout range};

      \nextgroupplot[
        height=0.135\columnwidth,
        ylabel={Draft sequence\\queue depth},
        ymin=0,ymax=6.8,
        ytick={0,3,6},
      ]
      \csTimeoutRange{8}{12}
      \addplot[cs step] table[x=shot,y=queue_depth,col sep=comma]
        {\caseStudyDataDir/12mp_normal_backlog.csv};

      \nextgroupplot[
        height=0.16\columnwidth,
        ylabel={Applied\\delay (ms)},
        ymin=0,ymax=1100,
        ytick={0,500,1000},
      ]
      \csTimeoutRange{8}{12}
      \addplot[cs delay] table[x=shot,y=delay_ms,col sep=comma]
        {\caseStudyDataDir/12mp_normal_delay.csv};

      \nextgroupplot[
        height=0.40\columnwidth,
        xlabel={Capture index within the 30-capture run},
        ylabel={\% of\\Capture Timeout},
        ymin=-11,ymax=80,
        ytick={0,20,40,60},
        yticklabels={0\%,20\%,40\%,60\%},
        legend style={
          at={(0.995,0.995)},anchor=north east,
          font=\tiny,draw=none,
          fill=white,fill opacity=0.75,text opacity=1,
          inner xsep=2pt,inner ysep=1pt,
          legend columns=2,
          /tikz/every even column/.append style={column sep=4pt},
        },
        legend cell align=left,
      ]
      \csTimeoutRange{8}{12}
      \addplot[cs floor,forget plot] coordinates {(0.5,0) (30.5,0)};
      \node[anchor=north east,font=\tiny,text=red!70!black,inner sep=1.5pt]
        at ([xshift=-1pt,yshift=-1pt]axis cs:30.5,0)
        {Capture Timeout (Slack \(=0\%\))};
      \addplot[cs backlog] table[x=shot,y expr={\thisrow{backlog_ms}/70},col sep=comma]
        {\caseStudyDataDir/12mp_normal_backlog.csv};
      \addlegendentry{Backlog}
      \addplot[cs slack] table[x=shot,y expr={\thisrow{margin_ms}/70},col sep=comma]
        {\caseStudyDataDir/12mp_normal_margin.csv};
      \addlegendentry{Slack}
      \node[anchor=east,font=\tiny,inner sep=1.5pt] at (axis cs:7.7,2.5)
        {2.1\%};
      \node[anchor=west,font=\tiny,inner sep=1.5pt] at (axis cs:24.4,5.1)
        {5.1\%};
    \end{groupplot}

    \matrix[
      matrix of nodes,
      nodes={font=\tiny,anchor=west,inner ysep=0.4pt},
      column sep=2pt,row sep=0.6pt,anchor=south east,
    ] at ([xshift=5pt,yshift=-2pt]group c1r1.north east) {
      \node[rectangle,draw=black,fill=black!72,line width=0.4pt,
            minimum size=2.7pt,inner sep=0pt] at (0.12,0) {}; & stage executed \\
      \node[rectangle,draw=black!45,line width=0.4pt,
            minimum size=2.7pt,inner sep=0pt] at (0.12,0) {}; & stage skipped by admission \\
    };
  \end{tikzpicture}
  \caption{Control interaction in the case-study run.}
  \label{fig:casestudy_12mp}
\end{figure}

